

\section*{\SGAInlineRefTarget{sga3:iii:section:2}2. Infinitesimal extensions of a morphism of group schemes}
\addcontentsline{toc}{section}{2. Infinitesimal extensions of a morphism of group schemes}

Retain the notation of \S0. Let \(Y\) and \(X\) be two \(S\)-group
functors, and let \(M\) be the kernel of the group homomorphism
\(p_X:X\to X^+\). Thus there is an exact sequence of \(S\)-group
functors
\[
 1\longrightarrow M\longrightarrow X\xrightarrow{p_X}X^+.
\]
By the definition of \(X^+\), there are isomorphisms
\begin{align*}
 \operatorname{Hom}_S(Y,X^+)
  &\xrightarrow{\sim}\operatorname{Hom}_{S_{\mathscr J}}(Y_{\mathscr J},X_{\mathscr J}),\\
 \operatorname{Hom}_{S\text{-}\mathrm{gr}}(Y,X^+)
  &\xrightarrow{\sim}
    \operatorname{Hom}_{S_{\mathscr J}\text{-}\mathrm{gr}}
    (Y_{\mathscr J},X_{\mathscr J}).
\end{align*}
The map
\[
 \operatorname{Hom}_S(Y,p_X):
 \operatorname{Hom}_S(Y,X)\longrightarrow\operatorname{Hom}_S(Y,X^+)
\]
sends an \(S\)-morphism \(f:Y\to X\) to the \(S\)-morphism
\(f^+:Y\to X^+\) corresponding under these isomorphisms to the
\(S_{\mathscr J}\)-morphism
\(f_{\mathscr J}:Y_{\mathscr J}\to X_{\mathscr J}\) obtained from \(f\)
by base change. If \(M\) is commutative, \SGARefLink{sga3:iii:proposition:1-2-4}{Proposition 1.2.4} applies to
this situation.

\medskip
\noindent\SGAInlineRefTarget{sga3:iii:paragraph:2-0}\textbf{2.0.}
In what follows we shall consider the case in which \(Y\) is flat over
\(S\), and \(X\) is an \(S\)-group functor of one of the following two
types:
\begin{enumerate}
\item[\(a\)] \(X\) is an \(S\)-group scheme;
\item[\(b\)]\SGAInlineRefTarget{sga3:iii:paragraph:2-0:item:b} \(X=\underline{\operatorname{Aut}}_S(E)\), where \(E\) is an
  \(S\)-scheme of finite presentation over \(S\).
\end{enumerate}
Write \((\mathrm{Flat})/S\) for the category of \(S\)-schemes flat over
\(S\). In case \(a\), respectively \(b\), the \(S\)-group functor
\(M=\operatorname{Ker}(X\to X^+)\), its restriction \(L\) to
\((\mathrm{Flat})/S\), and the actions on \(M\) of the inner
automorphisms of \(X\) were calculated in \SGARefLink{sga3:iii:corollary:0-9}{0.9}, \SGARefLink{sga3:iii:remark:0-5}{0.5}, and \SGARefLink{sga3:iii:corollary:0-10}{0.10},
respectively in \SGARefLink{sga3:iii:proposition:0-11}{0.11}, \SGARefLink{sga3:iii:remark:0-14}{0.14}, and \SGARefLink{sga3:iii:corollary:0-12}{0.12}.

More explicitly, in case \(a\), let \(L_0\) be the \(S_0\)-functor in
commutative groups defined, for every \(S_0\)-scheme \(T_0\), by
\[
 \operatorname{Hom}_{S_0}(T_0,L_0)
 =\operatorname{Hom}_{\mathscr O_{T_0}}
 \left(
  \omega^1_{X_0/S_0}\otimes_{\mathscr O_{S_0}}\mathscr O_{T_0},
  \mathscr J\otimes_{\mathscr O_{S_0}}\mathscr O_{T_0}
 \right).
\]
The group \(X_0\) acts on this functor through its adjoint
representation on \(\omega^1_{X_0/S_0}\), and
\[
 L=\prod_{S_0/S}L_0;
 \qquad L(T)=L_0(T\times_S S_0)
\]
for every \(S\)-scheme \(T\).

In case \(b\), let \(\pi:E\to S\) be the structural morphism. Then \(L\)
is the abelian-group functor on \((\mathrm{Flat})/S\) defined by
\[
 \operatorname{Hom}_S(T,L)
 =\Gamma\!\left(
  T,
  \pi_*\operatorname{Hom}_{\mathscr O_E}
   (\Omega^1_{E/S},\mathscr J\mathscr O_E)
  \otimes_{\mathscr O_S}\mathscr O_T
 \right),
\]
on which \(X\), regarded as a functor on \((\mathrm{Flat})/S\), acts as
described in \SGARefLink{sga3:iii:corollary:0-12}{0.12}. In either case there is an exact sequence of group
functors on \((\mathrm{Flat})/S\):
\[
 \tag{E}\label{sga3:iii:paragraph:2-0:equation:e}
 1\longrightarrow L\longrightarrow X\longrightarrow X^+.
\]

Since \(Y\) is flat over \(S\), \SGARefLink{sga3:iii:lemma:1-1-1}{Lemma 1.1.1} shows that the groups
\(H^i(Y,M)\) depend only on the restriction \(L\) of \(M\) to
\((\mathrm{Flat})/S\). In case \(a\), where
\(L=\prod_{S_0/S}L_0\), \SGARefLink{sga3:iii:lemma:1-1-2}{Lemma 1.1.2} gives isomorphisms
\[
 H^i(Y,L)\simeq H^i(Y_0,L_0).
\]

Taking all this into account, \SGARefLink{sga3:iii:proposition:1-2-4}{Proposition 1.2.4} yields the following
result.\footnote{Editor's note: The definitions that appeared in the
original statement of \SGARefLink{sga3:iii:theorem:2-1}{Theorem 2.1} and occupied most of original page
113 have been moved above.}

\medskip
\noindent\SGAInlineRefTarget{sga3:iii:theorem:2-1}\textbf{Theorem 2.1.}
\textit{Let \(S\) be a scheme, and let \(\mathscr I\supset\mathscr J\)
be two quasi-coherent ideals satisfying
\(\mathscr I\mathscr J=0\), defining the closed subschemes \(S_0\) and
\(S_{\mathscr J}\). Let}
\begin{itemize}
\item[] \textit{\(X\) be an \(S\)-group functor of type \(a\) or \(b\),
  with \(L_0\) and \(L\) as above;}
\item[] \textit{\(Y\) be an \(S\)-group scheme flat over \(S\), and
  \(f_{\mathscr J}:Y_{\mathscr J}\to X_{\mathscr J}\) a morphism of
  \(S_{\mathscr J}\)-groups.}
\end{itemize}
\textit{Then:}
\begin{enumerate}
\item[\(i\)] \textit{The morphism \(f_{\mathscr J}\) lifts to a
  morphism of \(S\)-groups \(Y\to X\) if and only if the following two
  conditions hold:}
  \begin{enumerate}
  \item[\(i_1\)] \textit{\(f_{\mathscr J}\) lifts to a morphism of
    \(S\)-functors \(Y\to X\). By \SGARefLink{sga3:iii:proposition:1-2-4}{Proposition 1.2.4} this defines an
    action of \(Y\) on \(L\), independent of the chosen lift. In case
    \(a\), the resulting action of \(Y_0\) on \(L_0\) comes from
    \(f_0:Y_0\to X_0\) and the adjoint action of \(X_0\) on \(L_0\).}
  \item[\(i_2\)] \textit{The canonically defined obstruction
    \(c(f_{\mathscr J})\in H^2(Y,L)\) vanishes. In case \(a\), this group
    is isomorphic to \(H^2(Y_0,L_0)\).}
  \end{enumerate}
\item[\SGAInlineRefTarget{sga3:iii:theorem:2-1:item:ii}(ii)] \textit{If the conditions in \(i\) hold, the set
  \(\mathcal E\) of morphisms of \(S\)-group functors \(Y\to X\)
  extending \(f_{\mathscr J}\) is principal homogeneous under
  \(Z^1(Y,L)\), isomorphic in case \(a\) to \(Z^1(Y_0,L_0)\). Moreover,
  the quotient of \(\mathcal E\) by the action of inner automorphisms
  of \(X\) defined by sections over \(S\) that induce the unit section
  of \(X_{\mathscr J}\) over \(S_{\mathscr J}\) is principal
  homogeneous under \(H^1(Y,L)\), isomorphic in case \(a\) to
  \(H^1(Y_0,L_0)\).}
\item[\SGAInlineRefTarget{sga3:iii:theorem:2-1:item:iii}(iii)] \textit{If \(f:Y\to X\) is a morphism of \(S\)-group
  functors extending \(f_{\mathscr J}\), the set of morphisms \(Y\to X\)
  obtained from \(f\) by those inner automorphisms is isomorphic to}
  \[
   \Gamma(L)/H^0(Y,L),
  \]
  \textit{and in case \(a\) to
  \(\Gamma(L_0)/H^0(Y_0,L_0)\).}
\end{enumerate}

\medskip
\noindent\SGAInlineRefTarget{sga3:iii:remark:2-1-1}\textbf{Remark 2.1.1.}\footnote{Editor's note: This remark,
analogous to \SGARefLink{sga3:iii:remark:4-5-1}{4.5.1}, has been added to introduce the notation
\(d(f,f')\) and \(\bar d(f,f')\), used in \SGARefLink{sga3:iii:paragraph:4-38}{4.38}. Consequently, the
part of \SGARefLink{sga3:iii:theorem:2-1:item:ii}{2.1(ii)} concerning \(\mathcal E\) itself was also added.}
If \(f,f':Y\to X\) are morphisms of \(S\)-group functors extending
\(f_{\mathscr J}\), one obtains a cocycle
\[
 d(f,f')\in Z^1(Y,L)
 \quad\bigl(\simeq Z^1(Y_0,L_0)\text{ in case (a)}\bigr)
\]
such that
\begin{equation}
 f'=d(f,f')\cdot f.
 \tag{*}\label{sga3:iii:remark:2-1-1:equation:star}
\end{equation}
Write \(\bar d(f,f')\) for the image of \(d(f,f')\) in
\(H^1(Y,L)\), or equivalently in \(H^1(Y_0,L_0)\) in case \(a\).

\medskip
\noindent\SGAInlineRefTarget{sga3:iii:remark:2-2}\textbf{Remark 2.2.}
Keep the preceding notation; in particular, \(Y\) is flat over \(S\).
In case \(b\), formula (\SGARefLink{sga3:iii:remark:0-14:equation:0-14-1}{0.14.1}) identifies \(L\) with the restriction to
\((\mathrm{Flat})/S\) of
\[
 \mathbf W\!\left(
  \pi_*\operatorname{Hom}_{\mathscr O_E}
  (\Omega^1_{E/S},\mathscr J\mathscr O_E)
 \right).
\]
\SGAInlineRefTarget{sga3:iii:remark:2-2:item:a}In case \(a\), suppose in addition that \(X\) is locally of finite
presentation over \(S\). Formula (\SGARefLink{sga3:iii:remark:0-6:equation:0-6-1}{0.6.1}) then identifies \(L\) with the
restriction to \((\mathrm{Flat})/S\) of
\[
 \prod_{S_0/S}\mathbf W\!\left(
  \operatorname{Hom}_{\mathscr O_{S_0}}
   (\omega^1_{X_0/S_0},\mathscr J)
 \right).
\]
In both cases the module to which \(\mathbf W\) is applied is
quasi-coherent, by EGA I, 9.1.1. Suppose further that \(Y\) is affine
over \(S\). By I, \SGARefLink{sga3:i:subsection:5-3}{5.3}, one obtains
\begin{align*}
 H^i(Y,L)=H^i(Y_0,L_0)
  &=H^i\!\left(Y_0,
    \operatorname{Hom}_{\mathscr O_{S_0}}
    (\omega^1_{X_0/S_0},\mathscr J)\right)
  &&\text{in case (a)},\\
 H^i(Y,L)
  &=H^i\!\left(Y,
    \pi_*\operatorname{Hom}_{\mathscr O_E}
    (\Omega^1_{E/S},\mathscr J\mathscr O_E)\right)
  &&\text{in case (b)}.
\end{align*}

\medskip
\noindent\SGAInlineRefTarget{sga3:iii:remark:2-3}\textbf{Remark 2.3.}
\begin{enumerate}
\item[(1)] By \SGARefLink{sga3:iii:proposition:0-16}{0.16} and \SGARefLink{sga3:iii:corollary:0-17}{0.17}, condition \(i_1\) is automatic when
  \(Y\) is affine and
  \[
   \begin{cases}
    X\text{ is smooth over }S, &\text{in case (a)},\\
    E\text{ is smooth and affine over }S, &\text{in case (b)}.
   \end{cases}
  \]
\item[(2)] Under these hypotheses, with \(Y\) still flat over \(S\)
  as in \SGARefLink{sga3:iii:paragraph:2-0}{2.0}, \SGARefLink{sga3:iii:remark:2-2:item:a}{formulas 2.2(a)} and (\SGARefLink{sga3:iii:remark:0-6:equation:0-6-2}{0.6.2}) give in case \(a\)
  \[
   H^i(Y,L)=H^i(Y_0,L_0)
   =H^i\!\left(Y_0,
    \operatorname{Lie}(X_0/S_0)
    \otimes_{\mathscr O_{S_0}}\mathscr J\right),
  \]
  and in case \(b\), by (\SGARefLink{sga3:iii:remark:0-14:equation:0-14-2}{0.14.2}), \SGARefLink{sga3:iii:lemma:1-1-2}{1.1.2}, and I, \SGARefLink{sga3:i:subsection:5-3}{5.3},
  \[
   H^i(Y,L)=H^i\!\left(Y_0,
    \pi_{0*}\operatorname{Hom}_{\mathscr O_{E_0}}
    (\Omega^1_{E_0/S_0},
     \mathscr J\otimes_{\mathscr O_{S_0}}\mathscr O_{E_0})\right).
  \]
  \footnote{Editor's note: The preceding material has been added; see
  Editor's Note 31 in \SGARefLink{sga3:iii:remark:0-14}{0.14}.}
\end{enumerate}

We now state several corollaries for a diagonalizable \(S\)-group
\(Y\) (I, \SGARefLink{sga3:i:subsection:4-4}{4.4}). If \(S\) is affine, then, by I, \SGARefLink{sga3:i:theorem:5-3-3}{5.3.3},
\(H^n(Y,F)=0\) for every \(n>0\) and every quasi-coherent
\(\mathscr O_S\)-module \(F\).

\medskip
\noindent\SGAInlineRefTarget{sga3:iii:corollary:2-4}\textbf{Corollary 2.4.}
\textit{Let \(S_0\hookrightarrow S\) be a closed subscheme defined by
a nilpotent ideal, and let \(Y\) be a diagonalizable \(S\)-group. Let}
\begin{enumerate}
\item[\(a\)] \textit{\(X\) be an \(S\)-group locally of finite
  presentation over \(S\), or}
\item[\(b\)]\SGAInlineRefTarget{sga3:iii:corollary:2-4:item:b} \textit{\(X=\underline{\operatorname{Aut}}_S(E)\), where
  \(E\) is an \(S\)-scheme locally of finite presentation.}
\end{enumerate}
\textit{If \(f:Y\to X\) is a morphism of \(S\)-groups whose base
change \(f_0:Y_0\to X_0\) is the unit morphism, then \(f\) is the unit
morphism.}

Indeed, the question is local on \(S\), and in case \(b\) also local on
\(E\). We may therefore suppose that \(S\) is affine and, in case \(b\),
that \(E\) is of finite presentation over \(S\). Introduce the closed
subschemes \(S_n\subset S\) defined by the powers of the ideal defining
\(S_0\). This reduces the assertion to the case in which \(S_0\) is
defined by a square-zero ideal, where it follows from \SGARefLink{sga3:iii:theorem:2-1}{Theorem 2.1} and
\SGARefLink{sga3:iii:remark:2-2}{Remark 2.2}.

If \(f_0\) is not required to be the unit morphism, one has the
following statement.

\medskip
\noindent\SGAInlineRefTarget{sga3:iii:corollary:2-5}\textbf{Corollary 2.5.}
\textit{Let \(S\) and \(S_0\) be as in \SGARefLink{sga3:iii:corollary:2-4}{2.4}, with \(S\) affine. Let
\(Y\) be a diagonalizable \(S\)-group, \(X\) an \(S\)-group functor,
and \(f_0:Y_0\to X_0\) a morphism of \(S_0\)-group functors.}
\begin{enumerate}
\item[\(i\)]\footnote{Editor's note: The original stated: ``Suppose one
of the following two properties holds: \(a\) \(X\) is an \(S\)-group
locally of finite presentation; \(b\)
\(X=\underline{\operatorname{Aut}}_S(E)\), where \(E\) is of finite
presentation over \(S\). Then \(f_0\) extends to a morphism of
\(S\)-groups \(Y\to X\) if and only if it extends to a morphism of
\(S\)-functors \(Y\to X\).'' It added that this followed step by step
from part \SGARefLink{sga3:iii:corollary:2-5:item:ii}{(ii)} of the theorem. That argument appears insufficient. For
example, if \(\mathscr I^3=0\) and
\(\mathscr J=\mathscr I^2\), and if \(f:Y\to X\) is an \(S\)-functor
morphism lifting \(f_0\), then \(f_0\) lifts to an
\(S_{\mathscr J}\)-group morphism
\(g_{\mathscr J}:Y_{\mathscr J}\to X_{\mathscr J}\); but must
\(g_{\mathscr J}\) then lift to an \(S\)-functor morphism \(g:Y\to X\)?
In any event, assertion \(i\) is not used later, where \(X\) is always
assumed smooth over \(S\).}
\item[\SGAInlineRefTarget{sga3:iii:corollary:2-5:item:ii}(ii)] \textit{Suppose either that \(X\) is a smooth \(S\)-group,
or that \(X=\underline{\operatorname{Aut}}_S(E)\) with \(E\) smooth
and affine over \(S\). Then \(f_0\) extends to a morphism of \(S\)-groups
\(Y\to X\), and any two extensions are conjugate by an inner
automorphism of \(X\) defined by a section over \(S\) inducing the unit
section of \(X_0\) over \(S_0\).}
\end{enumerate}

Introduce \(S_n\) as above.\footnote{Editor's note: What follows has
been slightly modified and expanded from the original.} For \SGARefLink{sga3:iii:corollary:2-5:item:ii}{(ii)}, a
scheme smooth over \(S\) is necessarily locally of finite presentation.
Thus, in case \(b\), a smooth affine \(E/S\) is of finite presentation,
so the hypothesis in \SGARefLink{sga3:iii:paragraph:2-0:item:b}{2.0(b)} is indeed satisfied.

Condition \(i_1\) of \SGARefLink{sga3:iii:theorem:2-1}{Theorem 2.1} is automatic by \SGARefLink{sga3:iii:proposition:0-16}{0.16} and \SGARefLink{sga3:iii:corollary:0-17}{0.17}.
Moreover, every section of \(X_{S_n}\) over \(S_n\) lifts to a section
of \(X_{S_{n+1}}\) over \(S_{n+1}\): this follows from smoothness in
case \(a\), and from \SGARefLink{sga3:iii:corollary:0-17}{0.17} in case \(b\). Consequently, if \(f\) and \(f'\)
are two lifts of \(f_0\), one may arrange step by step that
\(f_n=f'_n\) by lifting the inner automorphism supplied by
\SGARefLink{sga3:iii:theorem:2-1:item:ii}{Theorem 2.1(ii)}. This proves the assertion.

The same argument, together with \SGARefLink{sga3:iii:remark:2-3}{Remark 2.3}, gives:

\medskip
\noindent\SGAInlineRefTarget{sga3:iii:corollary:2-6}\textbf{Corollary 2.6.}
\textit{Let \(\mathscr I\) be a nilpotent ideal of \(S\) defining
\(S_0\), let \(Y\) be an affine \(S\)-group flat over \(S\), and let
\(X\) be a smooth \(S\)-group.}
\begin{enumerate}
\item[\(i\)] \textit{If, for every \(n\geq0\),}
 \[
  H^2\!\left(Y_0,
   \operatorname{Lie}(X_0/S_0)\otimes_{\mathscr O_{S_0}}
   \mathscr I^{n+1}/\mathscr I^{n+2}\right)=0,
 \]
 \textit{then every \(S_0\)-group morphism \(f_0:Y_0\to X_0\) lifts
 to an \(S\)-group morphism \(f:Y\to X\).}
\item[(ii)] \textit{If, for every \(n\geq0\),}
 \[
  H^1\!\left(Y_0,
   \operatorname{Lie}(X_0/S_0)\otimes_{\mathscr O_{S_0}}
   \mathscr I^{n+1}/\mathscr I^{n+2}\right)=0,
 \]
 \textit{then any two such lifts are conjugate by an inner
 automorphism of \(X\) defined by a section over \(S\) inducing the
 unit section of \(X_0\) over \(S_0\).}
\end{enumerate}

\medskip
\noindent\SGAInlineRefTarget{sga3:iii:lemma:2-7}\textbf{Lemma 2.7.}
\textit{Let \(S\) be affine, \(G\) an affine \(S\)-group,
\(\mathscr F\) a quasi-coherent \(\mathscr O_S\)-module, and
\(\mathscr L\) a locally free \(\mathscr O_S\)-module. Suppose \(G\)
acts on \(\mathscr F\) in the sense of Expos\'e I, hence acts on
\(\mathscr F\otimes_{\mathscr O_S}\mathscr L\). Let \(\Lambda\) be
the ring of \(S\), and let \(L\) be the projective \(\Lambda\)-module
defining \(\mathscr L\). Then there is a canonical isomorphism}
\[
 H^*(G,\mathscr F\otimes_{\mathscr O_S}\mathscr L)
 \simeq H^*(G,\mathscr F)\otimes_\Lambda L.
\]

Let \(\mathscr A\) be the \(\mathscr O_S\)-algebra
\(\mathscr A(G)\), and consider the complex \(\mathscr C\) of quasi-coherent
\(\mathscr O_S\)-modules
\[
 0\longrightarrow\mathscr F
 \longrightarrow\mathscr F\otimes\mathscr A
 \longrightarrow\mathscr F\otimes\mathscr A\otimes\mathscr A
 \longrightarrow\cdots.
\]
\footnote{Editor's note: The original proof has been expanded and
 simplified. It also invoked the isomorphisms
 \(H^n(\mathscr C\otimes\mathscr L)\simeq
 H^n(\mathscr C)\otimes\mathscr L\) and
\(H_n(\Gamma(-))\simeq\Gamma(H_n(-))\), where
\(H^n(\mathscr C)\) denotes the cohomology sheaf of the complex
\(\mathscr C\).}

By I, \SGARefLink{sga3:i:subsection:5-3}{5.3}, \(H^*(G,\mathscr F)\), respectively
\(H^*(G,\mathscr F\otimes\mathscr L)\), is the cohomology of
\(\Gamma(S,\mathscr C)\), respectively
\(\Gamma(S,\mathscr C\otimes\mathscr L)\). Since \(S\) is affine,
EGA I, 1.3.12 gives
\[
 \Gamma(S,\mathscr C\otimes\mathscr L)
 \simeq\Gamma(S,\mathscr C)\otimes_\Lambda L.
\]
Because \(L\) is projective, hence flat,
\[
 H^*(\Gamma(S,\mathscr C)\otimes_\Lambda L)
 \simeq H^*(\Gamma(S,\mathscr C))\otimes_\Lambda L,
\]
which proves the lemma.

Using the lemma, \SGARefLink{sga3:iii:corollary:2-6}{Corollary 2.6} becomes:

\medskip
\noindent\SGAInlineRefTarget{sga3:iii:corollary:2-8}\textbf{Corollary 2.8.}
\textit{Let \(S\) be affine and let \(\mathscr I\) be a nilpotent
ideal defining \(S_0\). Suppose every
\(\mathscr I^{n+1}/\mathscr I^{n+2}\) is locally free on \(S_0\). Let
\(Y\) be an affine \(S\)-group flat over \(S\), let \(X\) be a smooth
\(S\)-group, and let \(f_0:Y_0\to X_0\) be an \(S_0\)-group morphism.}
\begin{enumerate}
\item[\(i\)]\SGAInlineRefTarget{sga3:iii:corollary:2-8:item:i} \textit{If
\(H^2(Y_0,\operatorname{Lie}(X_0/S_0))=0\), then \(f_0\) lifts to an
\(S\)-group morphism \(Y\to X\).}
\item[\SGAInlineRefTarget{sga3:iii:corollary:2-8:item:ii}(ii)] \textit{If
\(H^1(Y_0,\operatorname{Lie}(X_0/S_0))=0\), then any two such lifts are
conjugate by an inner automorphism of \(X\) defined by a section over
\(S\) inducing the unit section of \(X_0\) over \(S_0\).}
\end{enumerate}

Taking \(Y=X\) gives:

\medskip
\noindent\SGAInlineRefTarget{sga3:iii:corollary:2-9}\textbf{Corollary 2.9.}
\textit{Let \(S\) and \(S_0\) be as above, and let \(X\) be a smooth
affine \(S\)-group.}
\begin{enumerate}
\item[\(i\)] \textit{If
\(H^1(X_0,\operatorname{Lie}(X_0/S_0))=0\), every endomorphism of \(X\)
over \(S\) that induces the identity on \(X_0\) is the inner
automorphism defined by a section over \(S\) inducing the unit section
of \(X_0\) over \(S_0\).}
\item[(ii)] \textit{If
\(H^2(X_0,\operatorname{Lie}(X_0/S_0))=0\), every
\(S_0\)-automorphism of \(X_0\) extends to an \(S\)-automorphism of
\(X\).}
\end{enumerate}

\medskip
\noindent\SGAInlineRefTarget{sga3:iii:remark:2-10}\textbf{Remark 2.10.}
The assertions involving \(H^1\) have converses by \SGARefLink{sga3:iii:theorem:2-1}{Theorem 2.1}. For
example, let \(S=I_{S_0}\) be the dual-number scheme over \(S_0\)
(II, \SGARefLink{sga3:iii:theorem:2-1}{2.1}), and let \(X\) be a flat \(S\)-group such that every
automorphism over \(S\) inducing the identity over \(S_0\) is an inner
automorphism defined by a section over \(S\) inducing the unit section
of \(X_0\). Then
\[
 H^1(X_0,\operatorname{Lie}(X_0/S_0))=0.
\]

\medskip
\noindent\SGAInlineRefTarget{sga3:iii:corollary:2-11}\textbf{Corollary 2.11.}
\textit{Let \(S,\mathscr I,\mathscr J\) be as in \SGARefLink{sga3:iii:theorem:2-1}{2.1}. Let \(Y\) be an
\(S\)-group scheme flat over \(S\), let \(X\) be an \(S\)-group scheme,
and let \(f:Y\to X\) be an \(S\)-group morphism. The set of morphisms
from \(Y\) to \(X\) obtained from \(f\) by conjugation with elements
\(x\in X(S)\) inducing the unit of \(X(S_{\mathscr J})\) is isomorphic
to the quotient}
\[
 \mathcal E=
 \operatorname{Hom}_{\mathscr O_{S_0}}(\omega^1_{X_0/S_0},\mathscr J)
 \big/
 \operatorname{Hom}_{\mathscr O_{S_0}}(\omega^1_{X_0/S_0},\mathscr J)
 ^{\operatorname{ad}(Y_0)}.
\]
\textit{Here the second group consists of those
\(\mathscr O_{S_0}\)-morphisms
\(\omega^1_{X_0/S_0}\to\mathscr J\) which, after every base change
\(S'\to S_0\), yield morphisms}
\[
 \omega^1_{X_0/S_0}\otimes_{\mathscr O_{S_0}}\mathscr O_{S'}
 \longrightarrow
 \mathscr J\otimes_{\mathscr O_{S_0}}\mathscr O_{S'}
\]
\textit{invariant under the action of \(Y_0(S')\) on the first
factor.}

By \SGARefLink{sga3:iii:theorem:2-1:item:iii}{Theorem 2.1(iii)}, this set is
\(\Gamma(L_0)/H^0(Y_0,L_0)\). Now
\[
 \Gamma(L_0)=
 \operatorname{Hom}_{\mathscr O_{S_0}}(\omega^1_{X_0/S_0},\mathscr J),
\]
and \(H^0(Y_0,L_0)\) is precisely
\(\Gamma(L_0)^{\operatorname{ad}(Y_0)}\) in the sense just stated.

\medskip
\noindent\SGAInlineRefTarget{sga3:iii:corollary:2-12}\textbf{Corollary 2.12.}
\textit{Under the hypotheses of \SGARefLink{sga3:iii:corollary:2-11}{2.11}, suppose in addition that
\(\omega^1_{X_0/S_0}\) is locally free of finite rank. Then}
\[
 \mathcal E\simeq
 \Gamma\!\left(S_0,
  \operatorname{Lie}(X_0/S_0)\otimes_{\mathscr O_{S_0}}\mathscr J
 \right)
 \big/
 H^0\!\left(Y_0,
  \operatorname{Lie}(X_0/S_0)\otimes_{\mathscr O_{S_0}}\mathscr J
 \right).
\]
Indeed,\footnote{Editor's note: The following sentence has been added.}
local freeness of finite rank gives
\[
 \operatorname{Hom}_{\mathscr O_{S_0}}
 (\omega^1_{X_0/S_0},\mathscr J)
 \simeq
 \operatorname{Lie}(X_0/S_0)\otimes_{\mathscr O_{S_0}}\mathscr J.
\]

\medskip
\noindent\SGAInlineRefTarget{sga3:iii:corollary:2-13}\textbf{Corollary 2.13.}
\textit{Suppose in addition that \(Y_0\) is diagonalizable. Then}
\[
 \mathcal E\simeq
 \Gamma\!\left(S_0,
  \operatorname{Lie}(X_0/S_0)\otimes_{\mathscr O_{S_0}}\mathscr J
 \right)
 \big/
 \Gamma\!\left(S_0,
  \operatorname{Lie}(X_0/S_0)^{\operatorname{ad}(Y_0)}
  \otimes_{\mathscr O_{S_0}}\mathscr J
 \right),
\]
\textit{where
\(\operatorname{Lie}(X_0/S_0)^{\operatorname{ad}(Y_0)}\) is the factor
of the decomposition in I, \SGARefLink{sga3:i:proposition:4-7-3}{4.7.3} corresponding to the zero
character of \(Y_0\).}

Indeed, if \(Y_0\simeq D_{S_0}(M)\), the cited result gives a direct-sum
decomposition
\[
 \operatorname{Lie}(X_0/S_0)
 =\operatorname{Lie}(X_0/S_0)_0
 \oplus
 \bigoplus_{\substack{m\in M\\m\ne0}}
 \operatorname{Lie}(X_0/S_0)_m.
\]
Tensoring with \(\mathscr J\) gives the analogous decomposition for
\(\operatorname{Lie}(X_0/S_0)\otimes\mathscr J\), whence
\[
 H^0\!\left(Y_0,
  \operatorname{Lie}(X_0/S_0)\otimes\mathscr J\right)
 \simeq
 \Gamma\!\left(S_0,
  \operatorname{Lie}(X_0/S_0)_0
  \otimes_{\mathscr O_{S_0}}\mathscr J\right).
\]

\medskip
\noindent\SGAInlineRefTarget{sga3:iii:corollary:2-14}\textbf{Corollary 2.14.}
\textit{Suppose in addition that \(S_0\) is affine. Then}
\[
 \mathcal E\simeq
 \Gamma\!\left(
  S_0,
  \left[
   \operatorname{Lie}(X_0/S_0)
   \big/
   \operatorname{Lie}(X_0/S_0)^{\operatorname{ad}(Y_0)}
  \right]
  \otimes_{\mathscr O_{S_0}}\mathscr J
 \right).
\]



